\chapter{Q-Learner}

\section{How it works}


\section{Implementation}
\subsection{State representation}
state aggregation -> reduce curse of dimensionality but brings perceptual aliasing
features :
\begin{itemize}
    \item empty\_count (number of empty cells goes from 0 to 15)
    \item max\_log (logarithmic value of the maximum tile (1 for 2, 2 for 4, 3 for 8, etc...)
    \item corner (1 if the maximum tile is in a corner, 0 otherwise)
    \item possible\_moves (number of possible moves, from 0 to 4)
\end{itemize}
let's say the max\_log is \leq 13 (score of 8192), then the state space is 16*13*2*5 = 2240 states.
We represent the state with different bases for each feature to make sure that each combination of features corresponds to a unique state index using the formula:
\[
\text{state} = \text{empty\_count} \cdot 140 + \text{max\_log} \cdot 10 + \text{possible\_moves} \cdot 2 + \text{corner}
\]
For example :
\begin{itemize}
    \item empty\_count = 10
    \item max\_log = 11
    \item possible\_moves = 2
    \item corner = 1
\end{itemize}
it gives us the state index 1515 (\[(14 \times 5 \times 2) + 1 \times (5 \times 2) + 2 \times 2 + 1
\])




\subsection{select action}
made it so it can only select from valid actions to avoid noise
\subsection{update q-table}
same as select action
la reward prend en compte la diff du nombre de cases vides, la diff du plus grand tile, si le max tile est dans un corner. Amélioration : compter la diff de nombre de merges possibles mais flemme en sah